% Conclusiones y recomendaciones de \nombreDayana (rúbrica del Entregable 3: 4 puntos, «del alumno»).
% Escribirlas en primera persona, con tus palabras: 2 o 3 párrafos alcanzan. Reemplaza la línea
% \pendiente{...} de abajo por tu texto. Preguntas guía, no hace falta responderlas todas:
%   - ¿Qué aprendiste sobre concurrencia con este trabajo (worker pool, barrera, Spin, speedup)?
%   - ¿Qué resultado te sorprendió y por qué?
%   - ¿Qué limitación ves en el trabajo y qué recomendarías para mejorarlo o escalarlo?
%   - En lo que hiciste tú (tu PR), ¿qué harías distinto?
% Para compilar y ver el PDF: cd tp/informe/tp && ./compilar.sh  (sale en build/main.pdf).
\subsubsection*{\nombreDayana}

Este trabajo me ayudó a entender que implementar concurrencia se trata más allá de crear
goroutines, es importante decidir qué información comparte cada una y cuándo deben sincronizarse. Evaluando la 
correspondencia entre la implementación en Go y el modelo Promela, comprendí mejor
el papel del \emph{worker pool}, los acumuladores privados y la barrera: los workers pueden procesar
los bloques en cualquier orden, pero el coordinador no debe actualizar los centroides hasta que
todos hayan terminado. También aprendí que un modelo debe conservar las decisiones esenciales del
programa sin copiar todos sus detalles. En este caso, fue posible estudiar la sincronización sin
representar los cálculos numéricos de K-means y comprobar que los puntos no se pierdan y que los
centroides no se escriban mientras están siendo leídos.

Me sorprendió que aumentar los workers no produzca una mejora proporcional: con cuatro se obtuvo
un speedup de \cifra{speedup-p4}$\times$ y con ocho de \cifra{speedup-p8}$\times$, pero con una menor
eficiencia. Esto demuestra que la coordinación y los núcleos físicos también limitan el rendimiento.
Como el modelo de Spin usa un dominio pequeño y no representa el cierre del pool, considero que sus
resultados respaldan el diseño, pero deben complementarse con las pruebas de la implementación en
Go. Para mejorar el trabajo, recomendaría modelar el cierre de las goroutines y realizar mediciones
con más datos y en equipos con más núcleos físicos. Si volviera a realizar mi aporte, incluiría
desde el inicio una tabla que relacione cada elemento de Go con su abstracción en Promela y con la
propiedad que permite verificar, dado que esa trazabilidad facilitaría comprender por qué el modelo sí
representa el comportamiento concurrente que se busca validar.
